%Перечень условных обозначений (при необходимости) оформляется в алфавитном порядке
\sectioncentered*{Перечень условных обозначений}
\addcontentsline{toc}{section}{Перечень условных обозначений}
\label{sec:reduction}

\begin{description}
    \item[API] - программный интерфейс приложения;
    \item[CSP] - задача удовлетворения ограничений;
    \item[JWT] - токен JSON Web Token;
    \item[LCW] - модель с длинным контекстным окном;
    \item[LLM] - большая языковая модель;
    \item[OAuth 2.0] - протокол авторизации внешних сервисов;
    \item[RAG] - генерация с дополнением извлеченным контекстом;
    \item[REST] - архитектурный стиль взаимодействия программных компонентов по сети;
    \item[СУБД] - система управления базами данных.
\end{description}
